\section{Problem formulation}
\label{sec:problem}

We formalize joint timetable planning and holding control as a bi-level Markov decision process (MDP) operating over a single bus route with time-varying passenger demand.
This section first describes the transit system model (\Cref{sec:transit_model}), then defines the upper- and lower-level MDPs (\Cref{sec:bilevel_mdp}), states the constraint specification (\Cref{sec:constraints}), and concludes with a discussion of the timescale separation that motivates our algorithmic design (\Cref{sec:timescale}).

%% ---------------------------------------------------------------
\subsection{Bus transit system model}
\label{sec:transit_model}

We consider a single bidirectional bus route consisting of $M$ stations indexed $\{1,\dots,M\}$.
Service is provided according to a timetable that dispatches a total of 264~trips per day, alternating between the two directions.
Passenger demand is drawn from time-varying origin--destination (OD) matrices that capture morning, midday, and evening peaks.
At any point in time, at most $\Nfleet^{\max}=25$ buses may be concurrently in service; each trip is characterized by a launch time and a travel direction.
The dispatch headway---the time between consecutive launches in the same direction---is the primary planning lever, while per-station holding times constitute the real-time control lever.

%% ---------------------------------------------------------------
\subsection{Bi-level MDP formulation}
\label{sec:bilevel_mdp}

We decompose the joint optimization into two coupled MDPs that share a common environment but differ in observation granularity, action semantics, and decision frequency.

\subsubsection{Upper-level MDP (timetable planning)}
\label{sec:upper_mdp}

The upper-level policy $\piU$ determines dispatch headways.
It is triggered at each dispatch event, yielding approximately 132 decision epochs per episode (one full service day).

\paragraph{State.}
The upper state encodes system-level summary statistics:
\begin{equation}
  \label{eq:upper_state}
  \sU = \Bigl[\,
    \tfrac{\text{hour}}{24},\;
    \tfrac{\text{demand}}{1000},\;
    \tfrac{\text{fleet\_on\_route}}{\Nfleet^{\max}},\;
    \tfrac{\text{prev\_headway}}{600},\;
    \text{unhealthy\_rate}
  \,\Bigr] \;\in\; \R^{5},
\end{equation}
where each component is normalized to facilitate learning.
The \emph{unhealthy\_rate} captures the fraction of active buses whose observed headway deviates from the target by more than a tolerance threshold.

\paragraph{Action.}
At each dispatch event the upper policy outputs a scalar offset that shifts the scheduled launch time of the dispatched bus:
\begin{equation}
  \label{eq:upper_action}
  \aU = \delta_t \;\in\; [-120,\, 120]\;\text{s}.
\end{equation}
The effective launch time of the bus is $t_{\text{base}} + \delta_t$, where $t_{\text{base}}$ is the bus's scheduled launch from a baseline timetable with $H = 360$\,s. The target headway $\htarget$ communicated to the lower level remains $360$\,s, so the upper policy adjusts \emph{when} a bus departs without changing the lower's regulation target. This formulation produces the same expressive power as a triple-headway action over peak/off-peak/transition periods (used by the search-based baselines below) while admitting per-trip granularity and online adaptation that the static triple cannot.

\paragraph{Reward.}
The upper reward is computed at the episode level from a measurement vector (defined in \Cref{sec:constraints}).
Because each upper action influences the system over an extended horizon through the downstream holding decisions it induces, assigning credit to individual dispatch events is nontrivial; we defer the reward construction to~\Cref{sec:method}.

\subsubsection{Lower-level MDP (holding control)}
\label{sec:lower_mdp}

The lower-level policy $\piL$ decides how long each bus should dwell beyond its minimum stop time at every station.
A decision is triggered at each station arrival for each active bus, producing a high volume of decision epochs throughout the day.

\paragraph{State.}
The lower state captures the local context of an individual bus:
\begin{equation}
  \label{eq:lower_state}
  \sL = \bigl[\,
    \text{bus\_id},\;
    \text{station\_id},\;
    \text{hour},\;
    \text{direction},\;
    h_{\mathrm{fwd}},\;
    h_{\mathrm{bwd}},\;
    \hat{d},\;
    \Delta h,\;
    v_1,\dots,v_K
  \,\bigr] \;\in\; \R^{8+K},
\end{equation}
where $h_{\mathrm{fwd}}$ and $h_{\mathrm{bwd}}$ denote the forward and backward headways, $\hat{d}$ is the estimated dwell time, $\Delta h$ is the deviation of the current headway from $\htarget$, and $v_1,\dots,v_K$ are recent inter-station speeds that encode local traffic conditions.

\paragraph{Action.}
The lower action is a scalar holding time:
\begin{equation}
  \label{eq:lower_action}
  \aL \;\in\; [0,\,60] \;\text{seconds}.
\end{equation}

\paragraph{Reward and cost.}
The immediate lower reward is a function of headway deviation from the target set by the upper level.
We define the per-step cost as
\begin{equation}
  \label{eq:lower_cost}
  c(\sL,\aL) = \min\!\Bigl(\frac{(h_{\mathrm{fwd}} - \htarget)^{2}}{(\htarget)^{2}},\;1\Bigr),
\end{equation}
which penalizes headway irregularity in a scale-invariant manner and is clipped at~1 for numerical stability.

%% ---------------------------------------------------------------
\subsection{Constraint specification}
\label{sec:constraints}

The bi-level MDP is subject to two soft constraints, both evaluated through a measurement vector collected at the end of each episode.

\paragraph{Measurement vector.}
At the conclusion of a service day, we compute:
\begin{equation}
  \label{eq:measurement}
  \zvec = \bigl[\,\text{avg\_wait\_min},\;\text{peak\_fleet},\;\text{bunching\_rate}\,\bigr].
\end{equation}

\paragraph{Soft constraint (fleet budget).}
The peak number of buses simultaneously on route should respect the per-episode fleet budget $\Nfleet \in [\Nfleet^{\min}, \Nfleet^{\max}]$ sampled at episode start. We enforce this softly via an overshoot$^2$ penalty in the upper reward:
\begin{equation}
  \label{eq:soft_fleet}
  \mathrm{overshoot}(\zvec, \Nfleet) = \max(0,\, \text{peak\_fleet} - \Nfleet)^2,
\end{equation}
weighted by an adaptive coefficient maintained by $\thetavec$-OGD (see~\Cref{sec:meas_proj}). To prevent unbounded violation under noisy demand, the environment imposes an outer ceiling $\Nfleet^{\max} + 3$ buses, which bounds but does not zero out the realised overshoot.

\paragraph{Soft constraint (headway uniformity).}
The expected per-step cost must remain below a threshold:
\begin{equation}
  \label{eq:soft_constraint}
  \E_{\piU,\piL}\!\bigl[\,c(\sL,\aL)\,\bigr] \;\leq\; \climit,
\end{equation}
which encourages even spacing of buses along the route.
Satisfaction of this constraint is enforced via Lagrangian relaxation within the lower-level policy (see \Cref{sec:method}).

%% ---------------------------------------------------------------
\subsection{Timescale separation}
\label{sec:timescale}

A defining characteristic of the bi-level formulation above is the pronounced separation of timescales between the two decision layers.
The upper policy $\piU$ acts at dispatch events, which occur roughly every five~minutes on average.
The lower policy $\piL$, by contrast, acts at every station arrival for every active bus---a frequency on the order of once every 30~seconds system-wide when multiple buses are in service.
Consequently, hundreds of lower-level decisions may intervene between two consecutive upper-level decisions, and the two layers never act simultaneously.

This asynchrony distinguishes our setting from the semi-Markov options framework~\citep{sutton1999between} and feudal reinforcement learning~\citep{dayan1992feudal,vezhnevets2017feudal}, which assume that (i)~the manager and worker share a common state space or a fixed abstraction thereof, and (ii)~the manager selects a subgoal that the worker executes to completion before the manager acts again.
Neither assumption holds here: the upper state $\sU\in\R^{5}$ summarizes fleet-level statistics whereas the lower state $\sL\in\R^{8+K}$ encodes bus-level kinematics, and both policies issue decisions driven by environment events rather than by a synchronous clock or a predefined option duration.

More concretely, a single upper action (a headway triple) remains in effect while many buses traverse multiple stations under independent lower-level control; the upper policy cannot pause and wait for all induced lower decisions to conclude.
Standard hierarchical RL methods that propagate value or advantage across levels therefore require modification to handle the fully asynchronous, event-driven interaction between $\piU$ and $\piL$.
We address this challenge through a goal-conditioned cross-level coupling, detailed in \Cref{sec:method}.
